% before: mt-dirichlet-character
The theorems concern Dirichlet characters of every modulus, the trivial and the imprimitive ones
included. Their statements depend on the value of a character at integers that share a factor with
the modulus, and on the primitive character that induces it. The definitions below fix these
conventions, which agree with Davenport's.

% before: mt-l-function
Each theorem is a statement about one function of $s$ for each character, defined on the whole
complex plane. Away from the pole of a trivial character it is the analytic continuation of the
Dirichlet series. At the pole a value is still assigned, and the hypotheses of the theorems are
arranged around that point.

% before: mt-value-at-one
Since $L(s,\chi)$ is defined at every point, a statement about its zeros either excludes the pole or
says something about the value there. The value at $s=1$ of the Riemann zeta function enters the
hypothesis of the reduction to primitive characters, and it explains the hypothesis $s\neq1$ in the
statement of the Riemann Hypothesis below.

% before: mt-euler-factors
An imprimitive character and the primitive character inducing it have $L$-functions that differ by
finitely many Euler factors. This comparison is the reduction to primitive characters in
Section~\ref{section-mt-imprimitive}, and it also locates the extra zeros of an imprimitive
$L$-function.

% before: mt-imprimitive-zeros
The $L$-function of an imprimitive character omits the Euler factors at the primes dividing its
modulus. Its zeros therefore include those of the finitely many factors in the comparison above, and
these decide on which regions a statement about all characters can be made.

% before: mt-trivial-zeros
The Riemann Hypothesis is a statement about the zeros of $\zeta$ other than its trivial zeros, so
its statement lists them. For other characters the list depends on the parity of the character and,
for an imprimitive character, on its Euler factors. The theorems for all characters are stated on a
strip that contains none of these zeros.

% before: mt-riemann-hypothesis
The first theorem is the Riemann Hypothesis for the zeta function of
Definition~\ref{mt-l-function}. Its hypotheses exclude the trivial zeros and the pole. Like the
theorems after it, it is deduced from the half-plane statement of Theorem~\ref{mt-half-plane}.

% before: mt-half-plane
The half-plane statement is the theorem that the rest of the book proves. It covers every Dirichlet
character, of every modulus, at every point to the right of the critical line except the pole of a
trivial character. On the line $\operatorname{Re}s=1$ it contains the classical nonvanishing behind
the prime number theorem in arithmetic progressions.

% before: mt-strip
The Generalized Riemann Hypothesis for Dirichlet $L$-functions is stated for all characters on the
open strip $0<\operatorname{Re}s<1$. The strip contains no trivial zeros, and none of the zeros that
the Euler factors of an imprimitive character contribute on the imaginary axis. So the statement has
the same form for primitive and imprimitive characters.

% before: mt-strip-primitive
For a primitive character, the statement on the strip is the Generalized Riemann Hypothesis in its
usual form, because the other zeros are the trivial ones. It is the case of
Theorem~\ref{mt-strip} for primitive characters, and it follows from
Theorem~\ref{mt-half-plane} by the functional equation alone.

% before: mt-euler-factor-nonzero
The Euler factors that separate an imprimitive character from its primitive character vanish only
on the imaginary axis. To the right of that axis they are nonzero, so the comparison of
Lemma~\ref{mt-euler-factors} transfers zeros and nonvanishing between the two $L$-functions.

% before: mt-imprimitive
The later chapters prove the half-plane statement for primitive characters. The comparison of Euler
factors extends it to every character, with the pole of a trivial character excluded. The
hypothesis includes the character modulo $1$ at $s=1$, where it is the statement $\zeta(1)\neq0$
about the value of Remark~\ref{mt-value-at-one}.

% before: mt-functional-equation
A zero of a primitive $L$-function to the left of the critical line is reflected by the functional
equation to a zero of the conjugate $L$-function to the right of it. The equation is stated for the
completed function, and it includes the character modulo $1$, so that $\zeta$ needs no separate
treatment in the strip.

% before: mt-zeta-functional-equation
For the zeros of $\zeta$ with $\operatorname{Re}s<1/2$, those to the left of the strip included, the
reflection is used in a form that involves $\zeta$ itself. In this form the trivial zeros appear as
zeros of the cosine factor.

% before: mt-strip-primitive-of-half-plane
The reflection argument for a primitive character comes first. In the strip the Gamma factor neither
vanishes nor has a pole, so a zero of $L(s,\chi)$ is a zero of the completed function, and the
functional equation moves it across the critical line.

% before: mt-strip-of-half-plane
For an imprimitive character, a zero in the strip is a zero of the $L$-function of the primitive
character, because the Euler factors do not vanish there. The reflection argument then applies to
that primitive character.

% before: mt-riemann-of-half-plane
The Riemann Hypothesis concerns all zeros of $\zeta$, those with $\operatorname{Re}s\leq0$
included. The functional equation of Proposition~\ref{mt-zeta-functional-equation} reflects a zero
with $\operatorname{Re}s<1/2$ to a zero with $\operatorname{Re}s>1/2$, unless the cosine factor
vanishes, which happens only at the trivial zeros and at $s=0$. Only the half-plane statement for the character modulo $1$ is needed.

% before: mt-riemann-via-strip
A second deduction of the Riemann Hypothesis goes through the statement on the strip for $\zeta$.
Outside the strip it uses the classical nonvanishing of $\zeta$ on $\operatorname{Re}s\geq1$ and
the functional equation of the completed function.
